\begin{abstract}
Question answering over a PDF with retrieval-augmented generation is a common small project:
split the document into chunks, embed them, retrieve the most similar few and pass them to a
language model. Such apps are rarely measured. We took one, an open-source Streamlit app for
research papers, and evaluated its own code on the 1,005 questions of the QASPER development
set, using the arXiv PDFs of the 281 papers, Llama~3.1~8B as the answering model and Qwen3-32B
as a judge checked against blind hand labels. As shipped, a misplaced loop sent the first
question of every conversation to the model without the question or the document; the model
returned \result{N} empty answers out of 1,005. After the fix, \result{X\%} of answers were
judged correct, against \result{Y\%} when the model was given the annotators' own evidence
paragraphs. At equal prompt length (about 4,000 characters), 500-character chunks with
cross-encoder reranking put a gold evidence passage in the context for \SweepSmallRerankBudget{}
of \SweepQuestions{} questions, against \SweepBaseBudget{} for the app's 1,000-character
chunks; on answers this \result{raised correctness to Z\%}. Loading nine other papers alongside
the one a question is about lowered evidence retrieval from \MultidocZeroSingle{} to
\MultidocNineMulti{}; letting the user restrict the search to one paper restored it. All code,
per-question outputs and prompts are released.
\end{abstract}
